% GENERATED FILE. Edit canonical lessons, then run npm run generate:books.
% canonical-source-hash: fnv1a64:308fa7ad52817ce6
% canonical-lessons: MR-C476-bubul, MR-C476-nakpudi, MR-C476-hirdi, MR-C476-dadhi, MR-C476-misi

\chapter{The Face Up Close}
\label{ch:mr-the-face-up-close}
\hlchaptermodality{476}

\begin{chapteropening}
\textbf{By the end of this chapter:} \emph{I can name the eyeball, a nostril, the gums, a beard and a moustache.}

Complete the last lesson of chapter 476: I can name the eyeball, a nostril, the gums, a beard and a moustache.
\end{chapteropening}

\section[\texorpdfstring{\mr{बुबुळ} (bubuḷ)}{बुबुळ (bubuḷ)}]{\mr{बुबुळ} (bubuḷ) — the eyeball; the iris}

\label{lesson:MR-C476-bubul}



\begin{quote}
Before the new one: say the Marathi for unconcerned, then the Marathi for to be scared of.
\end{quote}

\subsection*{You'll want to know: \mr{बुबुळ}}
\textbf{\mr{बुबुळ}} — \emph{bubuḷ} — \textquotedblleft{}the eyeball; the iris\textquotedblright{}.

\begin{quote}
\mr{डोळ्याच्या} \mr{बुबुळावर} \mr{एक} \mr{लहान} \mr{डाग} \mr{आहे}.
\end{quote}

\begin{quote}
\emph{ḍoḷyācyā bubuḷāvar ek lahān ḍāg āhe.}
\end{quote}

\begin{quote}
There is a small spot on the eyeball.
\end{quote}

Neuter.

Neuter: \mr{डोळ्याचं} \mr{बुबुळ}.

\begin{cousinweb}
\textbf{\mr{बुबुळ}} (\emph{bubuḷ}) is a Marathi noun whose origin is not settled.
\end{cousinweb}

\begin{grammarlens}[title={What you've built}]
One more word for describing what you lived and felt.
\end{grammarlens}

\subsection*{Your turn}
\begin{itemize}
\raggedright
  \item \emph{Say it:} \emph{bubuḷ}
  \item \emph{Say it:} \emph{bubuḷ}, once more
  \item \emph{Say it:} \emph{ḍoḷyācyā bubuḷāvar ek lahān ḍāg āhe.} — There is a small spot on the eyeball.
  \item \emph{Recall:} say the Marathi for unconcerned, then the Marathi for to be scared of, then say \emph{bubuḷ} again
\end{itemize}

\begin{tcolorbox}[breakable,colback=teal!4,colframe=teal!35!black,title={Before you move on}]
What does \mr{बुबुळ} mean? (\textquotedblleft{}The eyeball; the iris\textquotedblright{}.) Say it once more.
\end{tcolorbox}
\section[\texorpdfstring{\mr{नाकपुडी} (nākpuḍī)}{नाकपुडी (nākpuḍī)}]{\mr{नाकपुडी} (nākpuḍī) — a nostril}

\label{lesson:MR-C476-nakpudi}



\begin{quote}
Before the new one: say the Marathi for to be scared of, then the Marathi for the eyeball.
\end{quote}

\subsection*{You'll want to know: \mr{नाकपुडी}}
\textbf{\mr{नाकपुडी}} — \emph{nākpuḍī} — \textquotedblleft{}a nostril\textquotedblright{}.

\begin{quote}
\mr{सर्दीमुळे} \mr{एक} \mr{नाकपुडी} \mr{बंद} \mr{झाली}.
\end{quote}

\begin{quote}
\emph{sardīmuḷe ek nākpuḍī band jhālī.}
\end{quote}

\begin{quote}
With the cold, one nostril got blocked.
\end{quote}

Feminine. Plural: \mr{नाकपुड्या}, \emph{nākpuḍyā}.

\mr{नाक} + \mr{पुडी}, a little packet.

\begin{cousinweb}
\textbf{\mr{नाकपुडी}} (\emph{nākpuḍī}) is a Marathi compound built from \mr{नाक}, \textquotedblleft{}nose\textquotedblright{}, which CDIAL lists under \emph{*nakka} (CDIAL 6909), and \mr{पुडी}, \textquotedblleft{}a little packet\textquotedblright{}.
\end{cousinweb}

\begin{grammarlens}[title={What you've built}]
One more word for describing what you lived and felt.
\end{grammarlens}

\subsection*{Your turn}
\begin{itemize}
\raggedright
  \item \emph{Say it:} \emph{nākpuḍī}
  \item \emph{Say it:} \emph{nākpuḍī}, once more
  \item \emph{Say it:} \emph{sardīmuḷe ek nākpuḍī band jhālī.} — With the cold, one nostril got blocked.
  \item \emph{Recall:} say the Marathi for to be scared of, then the Marathi for the eyeball, then say \emph{nākpuḍī} again
\end{itemize}

\begin{tcolorbox}[breakable,colback=teal!4,colframe=teal!35!black,title={Before you move on}]
What does \mr{नाकपुडी} mean? (\textquotedblleft{}A nostril\textquotedblright{}.) Say it once more.
\end{tcolorbox}
\section[\texorpdfstring{\mr{हिरडी} (hirḍī)}{हिरडी (hirḍī)}]{\mr{हिरडी} (hirḍī) — a gum (of the mouth)}

\label{lesson:MR-C476-hirdi}



\begin{quote}
Before the new one: say the Marathi for the eyeball, then the Marathi for a nostril.
\end{quote}

\subsection*{You'll want to know: \mr{हिरडी}}
\textbf{\mr{हिरडी}} — \emph{hirḍī} — \textquotedblleft{}a gum (of the mouth)\textquotedblright{}.

\begin{quote}
\mr{दंतवैद्याने} \mr{माझ्या} \mr{हिरड्या} \mr{तपासल्या}.
\end{quote}

\begin{quote}
\emph{dantvaidyāne mājhyā hirḍyā tapāslyā.}
\end{quote}

\begin{quote}
The dentist checked my gums.
\end{quote}

Feminine. Plural: \mr{हिरड्या}, \emph{hirḍyā}.

Usually plural: \mr{हिरड्या}, the gums.

\begin{cousinweb}
\textbf{\mr{हिरडी}} (\emph{hirḍī}) is a Marathi noun whose origin is not settled.
\end{cousinweb}

\begin{grammarlens}[title={What you've built}]
One more word for describing what you lived and felt.
\end{grammarlens}

\subsection*{Your turn}
\begin{itemize}
\raggedright
  \item \emph{Say it:} \emph{hirḍī}
  \item \emph{Say it:} \emph{hirḍī}, once more
  \item \emph{Say it:} \emph{dantvaidyāne mājhyā hirḍyā tapāslyā.} — The dentist checked my gums.
  \item \emph{Recall:} say the Marathi for the eyeball, then the Marathi for a nostril, then say \emph{hirḍī} again
\end{itemize}

\begin{tcolorbox}[breakable,colback=teal!4,colframe=teal!35!black,title={Before you move on}]
What does \mr{हिरडी} mean? (\textquotedblleft{}A gum (of the mouth)\textquotedblright{}.) Say it once more.
\end{tcolorbox}
\section[\texorpdfstring{\mr{दाढी} (dāḍhī)}{दाढी (dāḍhī)}]{\mr{दाढी} (dāḍhī) — a beard}

\label{lesson:MR-C476-dadhi}



\begin{quote}
Before the new one: say the Marathi for a nostril, then the Marathi for a gum.
\end{quote}

\subsection*{You'll want to know: \mr{दाढी}}
\textbf{\mr{दाढी}} — \emph{dāḍhī} — \textquotedblleft{}a beard\textquotedblright{}.

\begin{quote}
\mr{आजोबांची} \mr{पांढरी} \mr{दाढी} \mr{लांब} \mr{आहे}.
\end{quote}

\begin{quote}
\emph{ājobāñcī pāṇḍhrī dāḍhī lāmb āhe.}
\end{quote}

\begin{quote}
Grandfather's white beard is long.
\end{quote}

Feminine.

\mr{दाढी} \mr{करणे}, to shave; \mr{दाढी} \mr{वाढवणे}, to grow a beard.

\begin{cousinweb}
\textbf{\mr{दाढी}} (\emph{dāḍhī}) is inherited: CDIAL lists Marathi \emph{dāḍhī}, \textquotedblleft{}fang, tusk, beard\textquotedblright{}, under \emph{*dāṁṣṭra} (CDIAL 6250).
\end{cousinweb}

\begin{grammarlens}[title={What you've built}]
One more word for describing what you lived and felt.
\end{grammarlens}

\subsection*{Your turn}
\begin{itemize}
\raggedright
  \item \emph{Say it:} \emph{dāḍhī}
  \item \emph{Say it:} \emph{dāḍhī}, once more
  \item \emph{Say it:} \emph{ājobāñcī pāṇḍhrī dāḍhī lāmb āhe.} — Grandfather's white beard is long.
  \item \emph{Recall:} say the Marathi for a nostril, then the Marathi for a gum, then say \emph{dāḍhī} again
\end{itemize}

\begin{tcolorbox}[breakable,colback=teal!4,colframe=teal!35!black,title={Before you move on}]
What does \mr{दाढी} mean? (\textquotedblleft{}A beard\textquotedblright{}.) Say it once more.
\end{tcolorbox}
\section[\texorpdfstring{\mr{मिशी} (miśī)}{मिशी (miśī)}]{\mr{मिशी} (miśī) — a moustache}

\label{lesson:MR-C476-misi}



\begin{quote}
Before the new one: say the Marathi for a gum, then the Marathi for a beard.
\end{quote}

\subsection*{You'll want to know: \mr{मिशी}}
\textbf{\mr{मिशी}} — \emph{miśī} — \textquotedblleft{}a moustache\textquotedblright{}.

\begin{quote}
\mr{माझ्या} \mr{वडिलांना} \mr{जाड} \mr{मिशी} \mr{आहे}.
\end{quote}

\begin{quote}
\emph{mājhyā vaḍilānnā jāḍ miśī āhe.}
\end{quote}

\begin{quote}
My father has a thick moustache.
\end{quote}

Feminine. Plural: \mr{मिशा}, \emph{miśā}.

Often plural: \mr{मिशा}.

\begin{cousinweb}
\textbf{\mr{मिशी}} (\emph{miśī}) is a Marathi noun whose origin is not settled.
\end{cousinweb}

\begin{grammarlens}[title={What you've built}]
One more word for describing what you lived and felt.
\end{grammarlens}

\subsection*{Your turn}
\begin{itemize}
\raggedright
  \item \emph{Say it:} \emph{miśī}
  \item \emph{Say it:} \emph{miśī}, once more
  \item \emph{Say it:} \emph{mājhyā vaḍilānnā jāḍ miśī āhe.} — My father has a thick moustache.
  \item \emph{Recall:} say the Marathi for a gum, then the Marathi for a beard, then say \emph{miśī} again
\end{itemize}

\begin{tcolorbox}[breakable,colback=teal!4,colframe=teal!35!black,title={Before you move on}]
What does \mr{मिशी} mean? (\textquotedblleft{}A moustache\textquotedblright{}.) Say it once more.
\end{tcolorbox}
